% Transcription — fonds Alexandre Grothendieck, shelfmark 12, batch 9 (pages 161-167, the last seven of the folder).
% Established from the facsimile archives/batches/12/batch-09.pdf (a mirror of
% https://grothendieck.umontpellier.fr/12.pdf), per the skill
% transcribe-grothendieck. The batch file carries no cover sheet: PDF page N
% is page 160+N, as the Montpellier footer printed on each PDF page confirms,
% and the archivists' pencil numbers bottom-left agree with it on every leaf
% (161 to 167). \page{} numbers that position.
%
% Pass: Opus 5.5 (claude-opus-5-5), 2026-09-23 — first pass, unchecked against the pages by a human.
%
% What the batch is.
%   161, 163, 165, 167  Typed versos in English, scanned upside down: pages
%           VII.20, VII.18, VII.17, VII.15 (descending, as versos of a run
%           written on their backs) of somebody else's typescript on
%           regular and unique factorisation domains, reflexive modules and
%           the Brauer group B(A) ⊂ B(K), with the colophon « Radcliffe
%           Institute for Independent Study, Wellesley College » on VII.20.
%           No mark in his hand on any of them. Absent, per the settled rule
%           on typed versos (hand.md §6). They are presumably not the
%           inventory's « tapuscrit ».
%   162, 164, 166  Three leaves in his hand, the end of the run on positive
%           bilinear forms on motives that batch 8 carries (p. 158 « Formes
%           bilinéaires positives sur les motifs », p. 160 the axioms
%           (i)-(iii) on a set Φ_X of « formes fondamentales »
%           φ = (φ_i), φ_i ∈ Hom(h^i(X) ⊗ h^i(X), Λ(-i)): stability under ⊕
%           and ⊗, and (iii) Tr(s_u u) ≥ 0 with equality iff u = 0, s_u the
%           transpose of u : X → Y relative to φ, ψ).
%     162   (iii) restated as (iii bis): on Hom(Λ(-i), h^{2i}(X)) ≃
%           Γ h^{2i}(X)(i) the form induced by φ is definite positive; a
%           marginal note says the equivalence comes from (ii), a special
%           case of (iii) and duality (morphisms X → Y as elements of
%           Γ H^{2i}(X⊠Y)(i)). Then two forms φ, ψ ∈ Φ_X, ψ̃ = φ̃ ∘ A: the
%           transposes u', u'' = u'A, Tr uu'A ≥ 0, so A lies in the positive
%           cone of the Rosati-type involutive algebra End(h^i(X)) — φ and ψ
%           « equivalent » as Riemannian structures.
%     164   « Plus généralement »: M an effective motive of degree i, embedded
%           λ : M → h^i(X), μ : M → h^i(Y); the induced φ_0, ψ_0 are definite
%           and equivalent, via Tr (μuλ')'(μuλ') = Tr λ'λ u'μ'μ u = Tr u'u
%           (λ'λ = μ'μ = id). Hence a class Ψ_M of Riemannian forms on every
%           effective motive. One diagram (redrawn in tikz-cd).
%     166   The analogous conditions (i)-(iii bis) for the Ψ_M; the twist
%           B(M(-1), M(-1)) ≃ B(M, M) and whether Ψ_{M(-1)} matches Ψ_M; a
%           boxed paragraph cancelled by diagonals; reduction via (ii) to
%           1 ∈ B(Λ(-1), Λ(-1)) lying in Ψ_{Λ(-1)}; hence the notion of
%           definite positive form extends to non-effective motives. The
%           folder ends here.
%
% Notation fixed for the batch, following batch 8's p. 160: \underline{h}^{i}(X)
% for his underlined h (math only); his script Φ (drawn like a gothic P) is
% \Phi_{X}; the class of forms on a motive, drawn as a large trident, is
% \Psi_{M}; the Tate object Λ(i), Λ(-1) is \Lambda; B(M, N) his bilinear forms;
% the external product X ⊠ Y is \boxtimes.
%
% Legibility. A fast draft throughout: the formulas carry, much of the
% connective prose does not, and the struck lines on 164 and the cancelled box
% on 166 are only partly legible. The slanted margin of 162 was re-rotated
% (magick -rotate 20) and reads nearly through.
\documentclass[11pt,a4paper]{article}
% Le préambule commun à toutes les transcriptions du fonds Grothendieck.
%
% Il tient en une idée : **l'incertitude se compose, elle ne se résout pas.**
% Une transcription qui lisse un mot douteux a détruit la seule chose pour
% laquelle on la faisait — un lecteur doit pouvoir savoir, à chaque endroit,
% ce qui était sur le feuillet et ce qui a été ajouté. Les six macros qui
% suivent sont donc l'appareil critique, et non de la décoration.
%
% Les mêmes commandes sont reconnues par `scripts/render.mjs`, qui produit la
% vue de lecture du site. Une macro ajoutée ici doit l'être aussi là-bas,
% faute de quoi le rendu échouera bruyamment — ce qui est voulu : un rendu
% silencieusement faux est pire qu'un rendu absent.

\ProvidesPackage{grothendieck}[2026/08/08 Transcriptions du fonds Grothendieck]

\usepackage{amsmath,amssymb,amsthm}
\usepackage{mathrsfs}
\usepackage{stmaryrd}
% Les diagrammes commutatifs sont partout dans ce fonds : c'est la langue
% même de la géométrie algébrique de Grothendieck, et une transcription qui
% les rendrait en prose ne transcrirait rien.
\usepackage{tikz-cd}
% Français par défaut — c'est la langue des feuillets — mais l'anglais est
% chargé aussi : la traduction et le résumé sont des documents anglais, et la
% typographie française (espaces avant les deux-points) y serait une faute.
% Ces documents basculent par \selectlanguage{english} en tête de corps.
\usepackage[english,french]{babel}
\usepackage{fontspec}
\usepackage{geometry}
\geometry{a4paper,margin=25mm,marginparwidth=28mm}
\usepackage{marginnote}
\usepackage{xcolor}
% ulem, not soul. `soul`'s \st and \ul are built on a character-by-character
% scan that XeLaTeX's font machinery breaks: the document fails with an
% undefined control sequence rather than degrading. ulem does the same two jobs
% and is Unicode-engine safe. `normalem` stops it hijacking \emph, which the
% transcriptions use for Grothendieck's own emphasis.
\usepackage[normalem]{ulem}
\usepackage{hyperref}
\hypersetup{colorlinks=true,linkcolor=black,urlcolor=black,pdfborder={0 0 0}}

% --- Métadonnées du lot -----------------------------------------------------
% Reprises telles quelles par le script de rendu, qui les affiche en tête de
% la vue de lecture. Elles fixent ce que le fichier prétend couvrir : sans
% elles, une transcription flotte sans point d'attache dans un fonds de seize
% mille feuillets.

\newcommand{\folder}[1]{\gdef\grothFolder{#1}}
\newcommand{\batch}[1]{\gdef\grothBatch{#1}}
\newcommand{\leaves}[2]{\gdef\grothFirst{#1}\gdef\grothLast{#2}}
\newcommand{\foldertitle}[1]{\gdef\grothFoldertitle{#1}}
\gdef\grothFolder{?}\gdef\grothBatch{?}\gdef\grothFirst{?}\gdef\grothLast{?}\gdef\grothFoldertitle{}

% --- Le résumé --------------------------------------------------------------
%
% Il ouvre la lecture modernisée, sous le titre « Résumé », et il est composé
% en petit. Ce n'est pas un ornement : c'est le seul passage du document qui
% ne soit pas des mathématiques, et un lecteur doit pouvoir le reconnaître
% avant de l'avoir lu — puis le sauter s'il connaît déjà le sujet, ou s'y
% arrêter s'il ne le connaît pas. Le filet qui le suit marque la reprise du
% corps.

\newenvironment{resume}
  {\small\setlength{\parskip}{0.45em}}
  {\par\medskip\hrule\medskip}

% --- Mots-clés ---------------------------------------------------------------
%
% En anglais, en clôture du résumé de la lecture modernisée : le vocabulaire
% moderne sous lequel chercher ce que les feuillets construisent. Le manifeste
% du site (`npm run manifest`) les extrait pour étiqueter la cote entière —
% cette ligne est la source unique des tags, il n'y a pas de fichier à côté.
\newcommand{\keywords}[1]{\par\smallskip\noindent
  {\footnotesize\textsc{Keywords} --- \textit{#1}}\par}

% --- L'appareil critique ----------------------------------------------------

% Le numéro de feuillet, en marge. C'est l'ancre qui permet de régler un
% désaccord sur une formule en regardant *un* feuillet plutôt qu'une chemise
% de six cents. Le site s'en sert aussi pour tourner les pages du fac-similé
% quand on fait défiler la transcription.
\newcommand{\leaf}[1]{%
  \marginnote{\footnotesize\color{gray!70}\textsf{#1}}%
  \label{leaf:#1}\ignorespaces}

% Illisible : on ne devine pas. Un mot inventé qui se lit comme les autres est
% le pire résultat possible.
\newcommand{\ill}{\textcolor{red!60!black}{[\,\ldots\,]}}

% Lecture douteuse : le mot est donné, et signalé comme incertain.
\newcommand{\uncertain}[1]{\uline{#1}}

% Ajout de l'éditeur — une lettre manquante, un mot sous-entendu.
\newcommand{\add}[1]{\textcolor{blue!50!black}{[#1]}}

% Biffé par Grothendieck lui-même. Ce qu'il a rayé fait partie du document :
% une rature dit souvent où la pensée a changé de direction.
\newcommand{\struck}[1]{\sout{#1}}

% Note du transcripteur, sur l'état matériel du feuillet.
\newcommand{\note}[1]{\par\smallskip{\small\color{gray!60!black}\textit{[#1]}}\par\smallskip}

% Note marginale de Grothendieck, distincte du corps du texte.
\newcommand{\marginal}[1]{\par\smallskip\noindent\hspace{1em}%
  {\small\color{gray!60!black}#1}\par\smallskip}

% --- Titre ------------------------------------------------------------------

\AtBeginDocument{%
  \begin{center}
    {\large\bfseries Fonds Alexandre Grothendieck --- cote n\textsuperscript{o}~\grothFolder}\\[2pt]
    {\itshape\grothFoldertitle}\\[4pt]
    {\small feuillets \grothFirst--\grothLast\ (lot \grothBatch)}\\[2pt]
    {\footnotesize\color{gray!60!black}Transcription établie d'après le fac-similé numérisé
     par l'Université de Montpellier.\\
     Les lectures incertaines sont soulignées, les illisibles notées [\,\ldots\,],
     les ajouts entre crochets.}
  \end{center}
  \vspace{1em}}

\endinput


\folder{12}
\batch{9}
\pages{161}{167}
\dating{1964-[vers 1971]}
\watermark{Édition de démonstration}
\foldertitle{Généralités (sous-groupes de Galois motiviques) (1964 ?) : notes manuscrites (s.d.), tapuscrit (s.d.)}

\begin{document}

\page{162}
\note{suite directe de la page 160 (lot 8), qui pose les conditions (i), (ii),
(iii) sur les ensembles $\Phi_{X}$ de formes fondamentales
$\varphi = (\varphi_{i})$,
$\varphi_{i} \in \operatorname{Hom}(\underline{h}^{i}(X) \otimes \underline{h}^{i}(X), \Lambda(-i))$}

On va exprimer l'axiome (iii) de façon équivalente, \uncertain{et} utilisant
$\Lambda(i)$.
\uncertain{en introduisant} $\chi = \varphi \boxtimes \psi$ sur $X \boxtimes Y$.
(On trouve alors \struck{sur $\operatorname{Hom}(\Lambda(-i), H^{2i}(X \boxtimes Y))$}
la forme
\note{« et utilisant $\Lambda(i)$ » est ajouté sous la fin de la première
ligne ; la parenthèse ouverte avant « On trouve » n'est pas refermée}

(iii bis) \emph{Sur} \struck{$H^{2i}(X)(i)$}
$\operatorname{Hom}(\Lambda(-i), \underline{h}^{2i}(X)) \simeq \Gamma\, \underline{h}^{2i}(X)(i)$,
\emph{la forme induite par la forme} $\varphi_{i}$ est \emph{définie}
\emph{positive}.
\note{il écrit bien $\varphi_{i}$, sur $\underline{h}^{2i}(X)$}

\marginal{l'équivalence de (iii) avec (iii bis) vient moyennant (ii) et le cas
\ill{} (iii$_{0}$), enfin \emph{dualité} pour exprimer les homomorphismes
$X \to Y$ en termes d'éléments de
$\Gamma H^{2i}(X \boxtimes Y)(i)$}
\note{note oblique de la marge gauche, en face de (iii bis) ; le petit
« 0 » sous « (iii) » se lit comme un indice, lecture douteuse ; le
« $\boxtimes$ » de la dernière ligne est surchargé}

Prenons $\varphi, \psi \in \Phi_{X}$, donc \struck{\ill{}}
$\widetilde{\psi} = \widetilde{\varphi} \circ A$. Alors pour tout
$u : X \to X$, induisant $u = u^{(i)} : \underline{h}^{i}(X) \to \underline{h}^{i}(X)$,
\struck{\ill{} transp.\ \ill{}}
soit $u'$ la transposée de $u$ relativement à $\varphi, \varphi$, et $u''$ la
transposée relatif à $\varphi, \psi$ $\to$ $u'' = u'A$, et on doit avoir
$\operatorname{Tr} u u'' \geqslant 0$ donc $\operatorname{Tr} u u' A \geqslant 0$.
\note{« donc $\widetilde{\psi} = \widetilde{\varphi} \circ A$ » est ajouté en
interligne et renvoyé par un trait après « $\Phi_{X}$ » ; le mot biffé
devant, peut-être $\varphi(x, x)$, ne se lit pas sûrement}

\emph{Cela exige que} $A$ \emph{est dans le cône positif} de l'algèbre
\uncertain{involutive} de Riemann $\operatorname{End}(\underline{h}^{i}(X))$
définie par $\varphi$, -- ce qui revient \uncertain{à dire que}, $\varphi$ et
$\psi$ sont des formes « équivalentes » \uncertain{au point de vue} des
structures Riemanniennes qu'elles définissent.
\note{« de Riemann » est ajouté au-dessus de la ligne ; un trait court sous
« involutive », qui peut être un soulignement ou une rature}

\page{164}
\struck{Soit maintenant} Plus généralement, \add{soit} $M$ un motif effectif,
\uncertain{homogène} de degré $i$. Choisissons un homomorphisme
\emph{injectif} $M \to \underline{h}^{i}(X)$ et considérons la forme induite
par une $\varphi \in \Phi_{X}$, soit $\varphi_{0}$. Choisissons de même
$M \to \underline{h}^{i}(Y)$, et $\psi \in \Phi_{Y}$, d'où $\psi_{0}$. Je dis
que $\varphi_{0}$ et $\psi_{0}$ sont (définies et) équivalentes.
\struck{\ill{} que \ill{} injective que} \struck{Tr}
\note{« Plus généralement » surcharge « Soit maintenant », biffé ; « (définies
et) » est ajouté au-dessus de la ligne et dans la marge}

\begin{tikzcd}[column sep=small, row sep=small]
  \underline{h}^{i}(X) \arrow[dr, bend left=20, "\lambda'"] & & & \underline{h}^{j}(Y) \arrow[dl, bend right=20, "\mu'"'] \\
  & M \arrow[ul, hook, bend left=20, "\lambda"] \arrow[r, "u"] & M \arrow[ur, hook', bend right=20, "\mu"'] & \\
  & \varphi_{0} & \psi_{0} &
\end{tikzcd}

\note{schéma redessiné d'après un croquis : $\varphi$ est écrit à gauche de
$\underline{h}^{i}(X)$, $\psi$ à droite de l'autre motif, $\varphi_{0}$ et
$\psi_{0}$ sous les deux $M$ ; le sens des flèches $\lambda$ et $\mu$ (crochet
côté $M$) est lu d'après le texte, $M \to \underline{h}^{i}(X)$ injectif ; un
symbole biffé précède $\lambda'$ ; l'exposant du motif de droite se lit
$j$ (ou $d$), là où le texte dit $\underline{h}^{i}(Y)$}

La première \ill{} \uncertain{viendrait} des \struck{\ill{}}
\struck{\ill{} que $\varphi$ est \ill{} \struck{définie}}
\struck{positive} … Riemanniennes …
La deuxième \struck{\ill{}} \struck{exprime} signifie que la forme
$\operatorname{Tr}(u'u)$ en $u \in \operatorname{End}(M)$ est définie, où $u'$
est défini par $\varphi_{0}$ et $\psi_{0}$.
\note{ces lignes, serrées dans la colonne de droite à côté du schéma, sont
en grande partie biffées ; « en $u \in \operatorname{End}(M)$ » est ajouté
au-dessus de « est définie »}

Or \struck{et} pour cela, il suffit de \struck{\ill{}} \uncertain{noter} que
la forme
\[
  \operatorname{Tr}\,(\mu u \lambda')'(\mu u \lambda')
  = \operatorname{Tr} \lambda u' \mu' \mu u \lambda'
  = \operatorname{Tr} \lambda' \lambda u' \mu' \mu u
\]
l'est, car $\mu'\mu = \mathrm{id}$ \add{et} $\lambda'\lambda = \mathrm{id}$,
donc cette forme est $\operatorname{Tr} u'u$, OK! Car \struck{\ill{}} donc,
\ill{}
\note{« $\lambda'\lambda = \mathrm{id}$ donc » est ajouté en interligne}

\emph{Pour les motifs effectifs}, on définit donc \struck{\ill{}}
\uncertain{des classes} d'équivalence de formes riemanniennes, \ill{} par
$\Psi_{M}$ (\ill{} \uncertain{notation} !). Les \uncertain{$\Psi_{M}$}
satisfont

\page{166}
Les conditions \struck{\ill{}} analogues à : (i) (ii) (iii) (iii bis)
\uncertain{commutativités}.
\note{le dernier mot est ajouté sous « (iii bis) », renvoyé par une flèche}

\struck{Lorsqu'on dispose d'un \ill{} de structure \ill{} par un \ill{}
$\Lambda(-1)$,} On a un isom
\[
  B(M(-1), M(-1)) \simeq B(M, M)
\]
et il y a lieu de vérifier la condition que $\Psi_{M(-1)}$ correspond à
$\Psi_{M}$. \struck{Cela \ill{}} doit pouvoir se faire \ill{} difficultés.
\struck{[\ill{} l'explicitation en termes de \ill{} directement, il faudrait
expliciter comment \ill{} $\Lambda(-1)$ en termes d'un \uncertain{dispositif}.]}
\note{le passage entre crochets est encadré et barré de trois diagonales ;
ses mots ne se lisent qu'en partie}

En fait, utilisant la propriété (ii), on \uncertain{devrait} \uncertain{voir}
qu'il suffit de vérifier que la forme canonique $1$ sur
\[
  1 \in B(\Lambda(-1), \Lambda(-1))
\]
est dans $\Psi_{\Lambda(-1)}$. \struck{Cela me semble} La vérification
dans le cas \uncertain{classique} \ill{} est triviale.
\note{« La vérification » est ajouté au-dessus de « Cela me semble »,
biffé ; le mot abrégé avant « est triviale » (« Alg. » ?) ne se lit pas}

Ceci permet d'étendre la notion de forme définie positive aux
motifs non nécessairement effectifs.
\note{« définie » est ajouté au-dessus de « positive »}

\end{document}
